\documentclass[10pt]{beamer}

\usepackage{verbatim}
%\usepackage[latin1]{inputenc}
\usepackage{amssymb,amsmath,amsfonts,url}
\usepackage{graphicx,color}
\usepackage{pgf,pgfarrows,pgfnodes,pgfautomata,pgfheaps,pgfshade}
\usepackage[utf8]{inputenc}
\usepackage[french]{babel}

\usepackage[all]{xy}
\usepackage{amstext}
\usepackage{xspace}
\usepackage{mdwtab}


\mode<article> {
  \usepackage{fullpage}
  \usepackage{hyperref}
  \setjobnamebeamerversion{advancedcourse.beamer}
}


\usepackage{listings}
\usepackage{xcolor}

\lstset{
	language=Python,
	basicstyle=\ttfamily\small,
	keywordstyle=\color{blue},
	commentstyle=\color{gray},
	stringstyle=\color{green!50!black},
	showstringspaces=false,
	frame=single,
	numbers=none
}
\lstset{
	language=Python,
	basicstyle=\ttfamily\small,
	keywordstyle=\color{blue},
	commentstyle=\color{gray},
	stringstyle=\color{green!50!black},
	showstringspaces=false,
	frame=single,
	numbers=none,
	literate=
	{é}{{\'e}}1
	{è}{{\`e}}1
	{ê}{{\^e}}1
	{ë}{{\"e}}1
	{à}{{\`a}}1
	{â}{{\^a}}1
	{ä}{{\"a}}1
	{ù}{{\`u}}1
	{û}{{\^u}}1
	{ü}{{\"u}}1
	{ô}{{\^o}}1
	{ö}{{\"o}}1
	{î}{{\^i}}1
	{ï}{{\"i}}1
	{ç}{{\c c}}1
}



\newcommand{\alertonce}[1]{\addtocounter{beamerpauses}{-1}\alert<+>{#1}}
\newcommand{\alerttwice}[1]{\addtocounter{beamerpauses}{-1}\alert<+-+(1)>{#1}}

%\newcommand{\subsubsubsection}[1]{\noindent\textbf{#1}}

\newcommand{\cfbox}[1]{\begin{center}\fbox{#1}\end{center}}
\newcommand{\cbox}[1]{\begin{center} #1 \end{center}}

%%%%%%%%%%%%%%%%%%%%%%%%%%%%%%%%%%%%%%%%%%%%%%%%%%%%%%%%%%%%%%%%%%%%%%%%%%%%%%%%%%%%%%%%%

\graphicspath{{Graphics/}}
%\input def_0.tex

%\input{StyleSlides.tex}
\mode<presentation> {
    \usepackage{beamerthemeBoadilla}

  \beamertemplatetransparentcovereddynamic
  \beamertemplateballitem
}


\title[IN100 Cours 1]{\textbf{Fondements informatiques I}\\
Cours 1: Variables et types}

\author[]{Sorina Ionica \texttt{sorina.ionica@uvsq.fr} \\~\\Bilal FAYE  \texttt{bilal.faye@uvsq.fr}}


\institute{}

\date{}


%\setbeamercolor{normal text}{fg=black}

\begin{document}


\frame{\titlepage}


\frame[containsverbatim]{
	\frametitle{Qu'est-ce qu'un algorithme ?}
	
	\begin{alertblock}{Un algorithme}
		Un algorithme est une méthode composée d'une suite finie d'étapes précises et ordonnées permettant de résoudre un problème.
		
		Pour décrire un algorithme, on précise :
		
		\begin{itemize}
			\item Entrées (Input) : les données dont on dispose au départ
			\item Traitement : les étapes à effectuer, dans un ordre précis
			\item Sortie (Output) : le résultat obtenu
		\end{itemize}
	\end{alertblock}
	
	\begin{exampleblock}{Exemple : calculer la moyenne de 3 notes}
		\small
		\begin{verbatim}
			Input : 12, 15 et 17
			Output : la moyenne des trois notes
			
			Additionner les trois notes : 12 + 15 + 17 = 44
			Diviser la somme par 3 : 44/3 = 14,67
			
			Renvoyer 14,67
		\end{verbatim}
	\end{exampleblock}
	
}

\frame[containsverbatim]{
	\frametitle{De l'algorithme au programme}
	
	Un \textbf{programme} est un ensemble d'instructions que l'ordinateur peut exécuter.
	
	\vspace{0.2cm}
	
	Pour écrire un programme, on utilise un \textcolor{blue}{langage de programmation}.
	
	\vspace{0.2cm}

\begin{exampleblock}{L'algorithme précédent en Python}
\begin{lstlisting}[language=Python]
# Input : Les données en entrée
note1 = 12
note2 = 15
note3 = 17
# Traitement : Additionner les trois notes
somme = note1 + note2 + note3
# Traitement : Diviser la somme par 3
moyenne = somme / 3
# Output : Afficher le résultat
print(moyenne)
\end{lstlisting}
\end{exampleblock}
\vspace{0.1cm}
$\rightarrow$ Un langage de programmation permet de traduire
un algorithme en instructions exécutables par l'ordinateur.

}


\frame{
	\frametitle{Pourquoi programmer ?}
	
	
	\begin{block}{Un algorithme devient utile lorsqu'il est exécuté par une machine.}
		Le programme traduit l'algorithme dans un \textbf{langage de programmation}
		que l'ordinateur peut comprendre et exécuter.
	\end{block}
	
	\vspace{0.35cm}
	
	\begin{columns}[T]
		\begin{column}{0.48\textwidth}
			\textbf{Exemples de programmes}
			
			\begin{itemize}
				\item Application web
				\item GPS et navigation
				\item Jeux vidéo
				\item Applications mobiles
				\item Réseaux sociaux
				\item Intelligence artificielle
			\end{itemize}
		\end{column}
		
		\begin{column}{0.48\textwidth}
			\textbf{Langages de programmation}
			
			\begin{itemize}
				\item Python
				\item Java
				\item C / C++
				\item JavaScript
				\item C\#
				\item Kotlin / Swift
			\end{itemize}
		\end{column}
	\end{columns}
	
	\vspace{0.3cm}
	
	\begin{exampleblock}{À retenir}
		\centering
		\textbf{Problème}
		$\;\longrightarrow\;$ 
		\textbf{Algorithme}
		$\;\longrightarrow\;$ 
		\textbf{Programme}
		$\;\longrightarrow\;$ 
		\textbf{Exécution par l'ordinateur}
	\end{exampleblock}
	
	
}


\frame{
	
	\begin{block}{Dans cette UE}
		Nous utiliserons \textbf{Python} pour écrire nos programmes.
	\end{block}
	
	\vspace{0.3cm}
	
	\begin{columns}[T]
		\begin{column}{0.48\textwidth}
			\textbf{Quelques repères}
			
			\begin{itemize}
				\item Créé par Guido van Rossum
				\item Première version en 1991
				\item Langage multiplateforme
			\end{itemize}
		\end{column}
		
		\begin{column}{0.48\textwidth}
			\textbf{Quelques caractéristiques}
			
			\begin{itemize}
				\item Gratuit et open-source
				\item Windows, Linux, macOS, Android...
				\item Langage interprété
			\end{itemize}
		\end{column}
	\end{columns}
	
	\vspace{0.3cm}
	
	\begin{exampleblock}{Langage interprété}
		Les instructions sont exécutées par un \textit{interpréteur},
		au fur et à mesure de l'exécution du programme.
	\end{exampleblock}
	
	
}

\begin{frame}[containsverbatim]
\frametitle{Les variables}

\begin{block}{Qu'est-ce qu'une variable ?}
	Une \textbf{variable} est un nom associé à une valeur dans un programme.
\end{block}
	
	\vspace{0.25cm}

\begin{exampleblock}{Exemple en Python}
\begin{lstlisting}[language=Python]
x = 2              # x prend la valeur 2
y = x              # y prend la valeur de x
print(x)           # affiche 2
print(y)           # affiche 2
y = 3              # y prend maintenant la valeur 3
print(x, y)        # affiche 2 3
\end{lstlisting}
\end{exampleblock}

\vspace{0.15cm}
\begin{center}
	 \begin{itemize}
	 	\item Ici, \texttt{x} et \texttt{y} sont des variables de type entier.
	 	\item Modifier \texttt{y} ne modifie pas \texttt{x}.
	 	\item \textit{Règle de style :} on met un espace après une virgule, mais pas avant.
	 \end{itemize}

\end{center}
\end{frame}
\begin{frame}[containsverbatim]
\frametitle{Échanger deux variables}
\begin{block}{Objectif}
Échanger les valeurs de deux variables sans perdre aucune valeur.
\end{block}

\vspace{0.25cm}

\begin{exampleblock}{Première méthode : incorrecte}
\begin{lstlisting}[language=Python]
x = 2
y = 3
x = y       # x prend la valeur de y
y = x       # y prend la valeur de x
print(x, y) # affiche 3 3
\end{lstlisting}
\end{exampleblock}
\vspace{0.15cm}
\begin{center}
\textbf{Problème :} la valeur initiale de \texttt{x} est perdue.
\end{center}
\end{frame}

\begin{frame}[containsverbatim]
\frametitle{Échanger deux variables}
\begin{block}{Objectif}
	Échanger les valeurs de deux variables sans perdre aucune valeur.
\end{block}
\vspace{0.25cm}
\begin{exampleblock}{Deuxième méthode : avec une variable temporaire}
\begin{lstlisting}[language=Python]
x = 2
y = 3
z = x       # sauvegarder la valeur de x
x = y       # x prend la valeur de y
y = z       # y récupère l'ancienne valeur de x
print(x, y) # affiche 3 2
\end{lstlisting}
\end{exampleblock}
\vspace{0.15cm}
\begin{center}
	\textbf{Idée :} utiliser \texttt{z} pour conserver temporairement une valeur.
\end{center}
\end{frame}

\begin{frame}[containsverbatim]
	\frametitle{Échanger deux variables}
	
	\begin{block}{Objectif}
		Échanger les valeurs de deux variables sans perdre aucune valeur.
	\end{block}
	\vspace{0.25cm}
	
\begin{exampleblock}{Troisième méthode : avec un tuple}
\begin{lstlisting}[language=Python]
x = 2
y = 3
x, y = y, x   # échange des deux valeurs
print(x, y)   # affiche 3 2
\end{lstlisting}
\end{exampleblock}
	\vspace{0.15cm}
	\begin{center}
		\textbf{En Python :} plusieurs valeurs peuvent être affectées simultanément.
	\end{center}
\end{frame}

\begin{frame}[containsverbatim]
	\frametitle{Règles de nommage des variables}
	\begin{block}{Comment choisir un nom de variable ?}
		Un nom de variable doit être \textbf{valide}, \textbf{lisible} et \textbf{explicite}.
	\end{block}
	\vspace{0.25cm}
	\begin{columns}[T]
		\begin{column}{0.48\textwidth}
			\textbf{Règles à respecter}
			
			\begin{itemize}
				\item Lettres : \texttt{a-z}, \texttt{A-Z}
				\item Chiffres : \texttt{0-9}
				\item Souligné : \texttt{\_}
				\item Pas d'espace ni de caractère accentué
				\item Le nom ne commence pas par un chiffre
			\end{itemize}
		\end{column}
		\begin{column}{0.48\textwidth}
			\textbf{Attention}
			
			\begin{itemize}
				\item Python est sensible à la casse
				\item \texttt{toto} $\neq$ \texttt{toTo}
				\item Ne pas utiliser de mot réservé
				\item Éviter les noms trop courts ou peu explicites
			\end{itemize}
		\end{column}
	\end{columns}
	
	\vspace{0.25cm}
	
\begin{exampleblock}{Exemples}
	\centering
	\textcolor{green!60!black}{\textbf{\checkmark Valide :}}
	\quad \texttt{age} \quad \texttt{nom\_etudiant} \quad \texttt{note1}
	
	
	\vspace{0.15cm}
	
	\textcolor{red}{\textbf{\texttimes Non valide :}}
	\quad \texttt{1note} \quad \texttt{nom etudiant} \quad \texttt{note-éleve}
	
	\vspace{0.15cm}
	
	\textcolor{red}{\textbf{\texttimes À éviter :}}
	\quad \texttt{print}
	\quad -- mot déjà utilisé par Python
	
	
\end{exampleblock}
\end{frame}

\begin{comment}
\frame[containsverbatim]{\frametitle{Régles de nommage des variables}
	
	Recommandations annexes
	\begin{itemize}
		\item sauf exceptions, donner un nom explicite aux variables tel \texttt{nb\_de\_vie}
		\item suivre cet exemple, cad séparer les mots par des  \texttt{\_}  sans majuscules
		\item éviter d'autres formes telles NbDeVie; au moins être cohérent dans tout le programme
	\end{itemize}
}
\end{comment}

\begin{frame}[containsverbatim]
\frametitle{Les types de données}

\begin{block}{Qu'est-ce qu'un type de donnée ?}
Le \textbf{type} d'une donnée détermine :
\begin{itemize}
\item les valeurs qu'elle peut prendre ;
\item les opérations que l'on peut lui appliquer.
\end{itemize}
\end{block}

\vspace{0.15cm}

\begin{exampleblock}{Les principaux types que nous allons manipuler au début}
\begin{lstlisting}[language=Python]
a = 5             # entier (int)
b = 3.14          # nombre flottant (float)
c = True          # booleen (bool)
d = "Bonjour"     # chaine de caracteres (str)
\end{lstlisting}
\end{exampleblock}

\vspace{0.1cm}

\begin{block}{En Python}
Le typage est \textbf{dynamique et implicite} :
\begin{itemize}
\item le type est déterminé automatiquement lors de l'affectation ;
\item une même variable peut changer de type au cours du programme.
\end{itemize}
\end{block}
\end{frame}
\begin{frame}[containsverbatim]
\frametitle{Les types de données}

\begin{block}{Connaître le type d'une donnée}
Pour connaître le type d'une donnée, on peut utiliser la fonction
\texttt{type()}.
\end{block}

\vspace{0.2cm}

\begin{exampleblock}{Exemple en Python}
\begin{lstlisting}[language=Python]
a = 5
b = 3.14
c = True
d = "Bonjour"
print(type(a))       # <class 'int'>
print(type(b))       # <class 'float'>
print(type(c))       # <class 'bool'>
print(type(d))       # <class 'str'>
\end{lstlisting}
\end{exampleblock}

\end{frame}


\begin{frame}[containsverbatim]
\frametitle{Les nombres}

En Python, on manipule principalement deux types de nombres :

\begin{columns}[T]
\begin{column}{0.48\textwidth}
\begin{block}{Entiers \texttt{int}}
Nombres sans partie décimale.

\textbf{Précision arbitraire} : Python peut représenter
des entiers très grands.
\end{block}
\end{column}

\begin{column}{0.48\textwidth}
\begin{block}{Flottants \texttt{float}}
Nombres avec une partie décimale.

Le \textbf{point} est utilisé comme séparateur décimal.
\end{block}
\end{column}
\end{columns}

\vspace{0.25cm}

\begin{exampleblock}{Exemples}
\begin{lstlisting}[language=Python]
x = 12          # int
y = 3.14        # float
z = 2e+20       # float : 2 x 10^20
\end{lstlisting}
\end{exampleblock}

\begin{center}
\textbf{Notation scientifique :} \texttt{2e+20} signifie $2 \times 10^{20}$.
\end{center}

\end{frame}
\begin{frame}[containsverbatim]
\frametitle{Opérateurs arithmétiques}

Sur les nombres, Python permet d'utiliser les opérateurs arithmétiques
usuels :

\begin{center}
\begin{tabular}{c|c|c}
\textbf{Opérateur} & \textbf{Opération} & \textbf{Exemple} \\ \hline
\texttt{+} & Addition & \texttt{5 + 2} \\
\texttt{-} & Soustraction & \texttt{5 - 2} \\
\texttt{*} & Multiplication & \texttt{5 * 2} \\
\texttt{/} & Division & \texttt{5 / 2} \\
\texttt{**} & Puissance & \texttt{5 ** 2} \\
\texttt{\%} & Reste & \texttt{5 \% 2}
\end{tabular}
\end{center}

\vspace{0.3cm}

\begin{exampleblock}{Exemple}
\begin{lstlisting}[language=Python]
a = 5 + 2       # 7
b = 5 * 2       # 10
c = 5 ** 2      # 25
d = 5 % 2       # 1
\end{lstlisting}
\end{exampleblock}

\end{frame}

\begin{frame}[containsverbatim]
\frametitle{Division et reste}

\begin{block}{Trois opérateurs à distinguer}
\begin{itemize}
\item \texttt{/} : division flottante
\item \texttt{//} : division entière
\item \texttt{\%} : reste de la division
\end{itemize}
\end{block}

\vspace{0.2cm}

\begin{exampleblock}{Exemple}
\begin{lstlisting}[language=Python]
print(11 / 2)    # 5.5
print(11 // 2)   # 5
print(11 % 2)    # 1
\end{lstlisting}
\end{exampleblock}

\vspace{0.15cm}

\begin{center}
$11 = 2 \times 5 + 1$

\vspace{0.1cm}

\textbf{Quotient : } 5
\qquad
\textbf{Reste : } 1
\end{center}

\end{frame}


\begin{frame}[containsverbatim]
\frametitle{Résultats des opérations numériques}

Le type du résultat dépend des types utilisés dans l'expression.

\begin{exampleblock}{Entiers et flottants}
\begin{lstlisting}[language=Python]
print(type(5 + 2))       # <class 'int'>
print(type(5 + 2.0))     # <class 'float'>
print(type(7 * 1.0))     # <class 'float'>
print(type(10 ** 2))     # <class 'int'>
\end{lstlisting}
\end{exampleblock}

\vspace{0.15cm}

\begin{alertblock}{À retenir}
La présence d'un \texttt{float} dans une opération numérique
peut conduire à un résultat de type \texttt{float}.
\end{alertblock}
\end{frame}
\begin{frame}[containsverbatim]
\frametitle{Opérateurs d'affectation augmentés}

Les opérateurs d'affectation augmentés permettent de modifier
une variable à partir de sa valeur actuelle.

\begin{exampleblock}{Ajouter une valeur}
\begin{lstlisting}[language=Python]
a = 2
b = 3
a += b       # équivaut à : a = a + b
print(a)     # 5
\end{lstlisting}
\end{exampleblock}

\vspace{0.2cm}

\begin{exampleblock}{Incrémenter une variable}
\begin{lstlisting}[language=Python]
a += 1       # équivaut à : a = a + 1
print(a)     # 6
\end{lstlisting}
\end{exampleblock}

\begin{center}
\texttt{-=} \qquad \texttt{*=} \qquad \texttt{/=} \qquad \texttt{//=} \qquad \texttt{\%=}
\end{center}

\end{frame}

\begin{frame}[containsverbatim]
\frametitle{Conversion de types}

Une donnée peut être convertie dans un autre type lorsque cela
est possible et pertinent.

\begin{block}{Deux fonctions de conversion}
\begin{itemize}
\item \texttt{int()} : conversion en entier
\item \texttt{float()} : conversion en flottant
\end{itemize}
\end{block}

\vspace{0.2cm}

\begin{exampleblock}{Exemple}
\begin{lstlisting}[language=Python]
x = 3
y = 3.14
a = float(x)       # 3 devient 3.0
b = int(y)         # 3.14 devient 3
print(type(a))     # <class 'float'>
print(type(b))     # <class 'int'>
\end{lstlisting}
\end{exampleblock}

\begin{alertblock}{Attention}
\texttt{int()} ne fait pas un arrondi : la partie décimale est supprimée.
\end{alertblock}
\end{frame}

\begin{frame}[containsverbatim]
\frametitle{La précision des nombres flottants}

Un nombre flottant est généralement représenté en machine
par une approximation en \textbf{base 2}.\newline


La précision étant limitée, certains nombres décimaux
ne peuvent pas être représentés exactement.

\vspace{0.2cm}

\begin{exampleblock}{Un résultat surprenant}
\begin{lstlisting}[language=Python]
print(0.1 + 0.2)
# 0.30000000000000004
print(0.1 + 0.2 == 0.3)
# False
\end{lstlisting}
\end{exampleblock}

\vspace{0.15cm}

\begin{alertblock}{À retenir}
Un \texttt{float} est souvent une \textbf{approximation} de la valeur
mathématique souhaitée.
\end{alertblock}

\end{frame}


\begin{frame}[containsverbatim]
\frametitle{Le type booléen}

\begin{block}{Une valeur booléenne}
Une donnée de type \texttt{bool} ne peut prendre que deux valeurs :
\[
\texttt{True} \qquad \text{ou} \qquad \texttt{False}
\]
\end{block}

\vspace{0.2cm}

Les booléens permettent notamment de représenter le résultat
d'une \textbf{condition}.

\begin{exampleblock}{Exemple}
\begin{lstlisting}[language=Python]
a = 5 < 2
print(a)           # False
print(type(a))     # <class 'bool'>
\end{lstlisting}
\end{exampleblock}

\vspace{0.15cm}

\begin{center}
\textbf{Une comparaison produit une valeur booléenne.}
\end{center}

\end{frame}


\begin{frame}[containsverbatim]
\frametitle{Les opérateurs de comparaison}

Les opérateurs de comparaison permettent de comparer deux valeurs.
Le résultat est toujours un booléen : \texttt{True} ou \texttt{False}.

\vspace{0.2cm}

\begin{center}
\begin{tabular}{c|c|c}
\textbf{Opérateur} & \textbf{Signification} & \textbf{Exemple} \\ \hline
\texttt{>}  & supérieur à & \texttt{5 > 2} \\
\texttt{<}  & inférieur à & \texttt{5 < 2} \\
\texttt{>=} & supérieur ou égal à & \texttt{5 >= 5} \\
\texttt{<=} & inférieur ou égal à & \texttt{5 <= 2} \\
\texttt{==} & égal à & \texttt{5 == 5} \\
\texttt{!=} & différent de & \texttt{5 != 2}
\end{tabular}
\end{center}

\vspace{0.25cm}

\begin{exampleblock}{Exemple en Python}
\begin{lstlisting}[language=Python]
age = 20
print(age >= 18)       # True
print(age == 20)       # True
print(age != 20)       # False
\end{lstlisting}
\end{exampleblock}

\end{frame}

\begin{frame}[containsverbatim]
\frametitle{Les opérateurs logiques}

Les opérateurs logiques permettent de \textbf{combiner ou modifier}
des valeurs booléennes.

\vspace{0.2cm}

\begin{columns}[T]
\begin{column}{0.31\textwidth}
\begin{block}{\texttt{not}}
Inverse une valeur.

\[
\texttt{not True = False}
\]
\[
\texttt{not False = True}
\]
\end{block}
\end{column}

\begin{column}{0.31\textwidth}
\begin{block}{\texttt{and}}
\texttt{True} seulement si
les deux valeurs sont vraies.

\[
\texttt{True and True}
=\texttt{True}
\]


\end{block}
\end{column}

\begin{column}{0.31\textwidth}
\begin{block}{\texttt{or}}
\texttt{True} si au moins
une valeur est vraie.

\[
\texttt{True or False}
=\texttt{True}
\]
\end{block}
\end{column}
\end{columns}

\vspace{0.25cm}

\begin{exampleblock}{Exemple}
\begin{lstlisting}[language=Python]
age = 20
majeur = age >= 18
print(majeur and True)      # True
print(majeur and False)     # False
print(not majeur)           # False
\end{lstlisting}
\end{exampleblock}

\end{frame}


\begin{frame}[containsverbatim]
\frametitle{Construire une expression logique}

Une expression logique peut combiner plusieurs comparaisons
avec \texttt{and}, \texttt{or} et \texttt{not}.

\vspace{0.2cm}

\begin{exampleblock}{Exemple}
\begin{lstlisting}[language=Python]
a = -0.5
p = (a <= 0)          # True
q = (a > 1)            # False
r = (a < 3)            # True
s = (a != 2)           # True
res = p or (q and r and s)
print(res)             # True
\end{lstlisting}
\end{exampleblock}

\vspace{0.15cm}

\begin{center}
\texttt{p} est vrai, donc \texttt{p or (...)} est vrai.
\end{center}

\vspace{0.1cm}

\begin{alertblock}{À retenir}
Une expression logique peut être construite à partir de
\textbf{comparaisons} et d'\textbf{opérateurs logiques}.
\end{alertblock}

\end{frame}

\begin{frame}[containsverbatim]
\frametitle{Les chaînes de caractères}

\begin{block}{Qu'est-ce qu'une chaîne de caractères ?}
Une \textbf{chaîne de caractères} (\texttt{str}) est une suite de caractères
écrite entre guillemets simples ou doubles.
\end{block}

\vspace{0.2cm}

\begin{exampleblock}{Exemples}
\begin{lstlisting}[language=Python]
s1 = "Bonjour"
s2 = 'Python'
s3 = "Bonjour Python"
print(s1)
print(type(s1))       # <class 'str'>
\end{lstlisting}
\end{exampleblock}

\vspace{0.1cm}

\begin{center}
\textbf{Simple ou double guillemets :} les deux sont possibles.
\end{center}

\end{frame}

\begin{frame}[containsverbatim]
\frametitle{Opérations sur les chaînes}

Les chaînes de caractères peuvent être manipulées avec plusieurs
opérations simples.

\begin{exampleblock}{Exemple}
\begin{lstlisting}[language=Python]
s = "Bonjour"
print(s + " Python")   # concaténation
print(s * 3)           # répétition
print(len(s))          # longueur : 7
\end{lstlisting}
\end{exampleblock}

\vspace{0.2cm}

\begin{block}{À retenir}
\begin{itemize}
\item \texttt{+} assemble deux chaînes ;
\item \texttt{*} répète une chaîne ;
\item \texttt{len()} donne le nombre de caractères.
\end{itemize}
\end{block}

\end{frame}

\begin{frame}[containsverbatim]
\frametitle{Accéder aux caractères d'une chaîne}

Une chaîne est une suite de caractères.
Chaque caractère possède une \textbf{position}, appelée indice.

\begin{exampleblock}{Exemple}
\begin{lstlisting}[language=Python]
s = "Python"
print(s[0])       # P
print(s[1])       # y
print(s[5])       # n
\end{lstlisting}
\end{exampleblock}

\vspace{0.15cm}

\begin{center}
\texttt{P \;\; y \;\; t \;\; h \;\; o \;\; n}

\vspace{0.1cm}

\texttt{0 \;\; 1 \;\; 2 \;\; 3 \;\; 4 \;\; 5}
\end{center}

\vspace{0.1cm}

\begin{alertblock}{Attention}
En Python, le premier indice est \textbf{0}.
\end{alertblock}

\end{frame}

\begin{frame}[containsverbatim]
\frametitle{Quelques fonctions pour les chaînes}

Une chaîne de caractères possède des \textbf{fonctions} permettant de la manipuler.

\begin{exampleblock}{Exemple}
\begin{lstlisting}[language=Python]
s = "bonjour"
print(s.upper())       # "BONJOUR"
print(s.capitalize())  # "Bonjour"
print(s.replace("jour", "soir"))
\end{lstlisting}
\end{exampleblock}

\vspace{0.15cm}

\begin{block}{Fonctions}
Une fonction peut être associée à une donnée.

\vspace{0.1cm}

La notation \texttt{s.upper()} signifie :
\textbf{appliquer la méthode \texttt{upper()} à \texttt{s}}.
\end{block}

\end{frame}


\begin{frame}[containsverbatim]
\frametitle{Lire une donnée au clavier}

La fonction \texttt{input()} permet de récupérer une donnée
saisie par l'utilisateur.

\begin{exampleblock}{Lire un texte}
\begin{lstlisting}[language=Python]
nom = input("Entrez votre nom : ")
print("Bonjour", nom)
\end{lstlisting}
\end{exampleblock}

\vspace{0.15cm}

\begin{alertblock}{Attention}
\texttt{input()} renvoie toujours une \textbf{chaîne de caractères}
(\texttt{str}), même si l'utilisateur saisit un nombre.
\end{alertblock}

\begin{exampleblock}{Lire un entier}
\begin{lstlisting}[language=Python]
age = input("Entrez votre âge : ")
age = int(age)       # conversion en entier
print(age + 1)
\end{lstlisting}
\end{exampleblock}

\end{frame}

\begin{frame}
	\frametitle{À retenir : les types de données}
	
	\begin{block}{Les types de base}
		\centering
		\texttt{int} \quad
		\texttt{float} \quad
		\texttt{bool} \quad
		\texttt{str}
	\end{block}
	
	\vspace{0.2cm}
	
	\begin{columns}[T]
		\begin{column}{0.48\textwidth}
			\textbf{En Python}
			\begin{itemize}
				\item Le type est déterminé automatiquement.
				\item Une variable peut changer de type.
				\item \texttt{type()} permet de connaître un type.
				\item \texttt{int()}, \texttt{float()} et \texttt{str()}
				permettent de convertir une donnée.
			\end{itemize}
		\end{column}
		
		\begin{column}{0.48\textwidth}
			\textbf{Opérations essentielles}
			\begin{itemize}
				\item Nombres : \texttt{+}, \texttt{-}, \texttt{*}, \texttt{/}, \texttt{//}, \texttt{\%}, \texttt{**}
				\item Booléens : \texttt{and}, \texttt{or}, \texttt{not}
				\item Chaînes : \texttt{+}, \texttt{*}, \texttt{len()}
				\item Comparaisons : \texttt{<}, \texttt{>}, \texttt{==}, \texttt{!=}, ...
			\end{itemize}
		\end{column}
	\end{columns}
	
	\vspace{0.25cm}
	
	\begin{exampleblock}{Idée essentielle}
		\centering
		\textbf{Le type d'une donnée détermine les opérations que l'on peut effectuer sur celle-ci.}
	\end{exampleblock}
	
\end{frame}

\end{document}
